\appendix
\section{Exact theorem schemas and packet checklist}
\label{app:schemas}

This appendix collects the core theorem shapes in a form suitable for direct
translation into Lean declarations.

\subsection{Boolean bounded-formula compiler}

For every bounded formula \(\theta(\bar x)\) formed from equations between
fixed primitive-recursive terms, propositional connectives, and bounded
quantifiers, construct a primitive-recursive Boolean
\(b_\theta(\bar x)\) and derivations
\[
 []\vdash[\PRAAx]
 \bigl(b_\theta(\bar x)=1\to\theta(\bar x)\bigr),
 \qquad
 []\vdash[\PRAAx]
 \bigl(\theta(\bar x)\to b_\theta(\bar x)=1\bigr).
\]
The construction is structural recursion on \(\theta\).  The bounded
quantifier cases use the fold-correctness derivations.

In the compiled Lean implementation, the source syntax should be a restricted
inductive type carrying its free-variable arity.  Compilation returns the
Boolean \(\PR\) symbol, the object-language formula, and both derivations.

\subsection{Proof-producing trace reflection}

Given fixed symbols \(T=(g,i,b,s,v,a)\), where \(g\) is the good-input bit,
\(i\) the initial state, \(b\) the bound, \(s\) the step, \(v\) the invariant,
and \(a\) the acceptance bit, let \(r=\code{PR.prec }i\ s\).  The input record
contains derivations of:
\begin{align*}
 g(u)=1&\to v(0,i(u),u)=1,\\
 g(u)=1\wedge v(j,z,u)=1&\to v(j+1,s(j,z,u),u)=1,\\
 g(u)=1\wedge v(b(u),z,u)=1&\to a(z,u)=1.
\end{align*}
The output is a derivation of
\[
 \forall u\;\bigl(g(u)=1\to a(r(b(u),u),u)=1\bigr).
\]
No semantic premise about the standard natural numbers appears in the theorem
signature; every premise is already an object-language derivation.

\subsection{Finite path compiler}

For fixed \(B,\Good,\Compat\), define the path checker
\(\mathsf{PathOK}(p,k,d,a,c)\).  Its fields assert:

\begin{enumerate}
\item the code \(c\) has length \(d+1\);
\item its first entry is \(a\);
\item every entry is below the candidate bound and good; and
\item every adjacent pair is compatible.
\end{enumerate}

The viability theorem is the pair of derivations
\begin{align*}
 \Viable(p,k,d,a)=1
 &\to
 \exists c<C(p,k,d)\;\mathsf{PathOK}(p,k,d,a,c)=1,\\
 \mathsf{PathOK}(p,k,d,a,c)=1
 &\to \Viable(p,k,d,a)=1,
\end{align*}
where \(C\) is a displayed primitive-recursive code bound.  The existential
witness in the first formula can be the least path code or the recursively
constructed least viable path.

\subsection{Shrinking viability}

The exact numerical conclusion is
\[
 \forall k\le\ell\quad
 d(\Obs(p,k,s(p,k)),\Obs(p,\ell,s(p,\ell)))
 \le\varepsilon(p,k)+
       \sum_{j=k}^{\ell-1}\delta(p,j).
\]
The assumptions are:

\begin{enumerate}
\item \(s(p,k)\) is the least member of the nonempty cloud
      \(\FullV(p,H(p,k),k,-)\);
\item \(H\) is nondecreasing and dominates its level;
\item the cloud diameter is at most \(\varepsilon(p,k)\);
\item compatible edges move observations by at most \(\delta(p,k)\); and
\item the coded distance satisfies symmetry and triangle inequality on the
      guarded domain used in the theorem.
\end{enumerate}

The regular-name corollary takes a separate tail modulus.  The theorem does not
infer summability of \(\delta\) or convergence of \(\varepsilon\).

\subsection{Finite-net transfer}

A net certificate has a finite point list \(N\), a coverage checker, and a
radius \(\rho\).  A modulus certificate has a bound \(\eta\) at radius
\(\rho\).  The embedded target is:
\[
 \mathsf{NetOK}(p,N,\rho)=1\wedge
 \mathsf{ModulusOK}(p,\rho,\eta)=1\wedge
 \All_{y\in N}[f(y)\ge L]=1
 \to
 \mathsf{GlobalLower}(p,f,L-\eta)=1.
\]
The meaning of \(\mathsf{GlobalLower}\) depends on the represented domain; its
proof must reduce an arbitrary challenge to one finite coverage witness.

\subsection{Height bound and list}

For a normalized point-code predicate \(P(p,x)\) and height \(h(x)\), define:
\begin{align*}
 \mathsf{BoundOK}(p,H)&:=\forall x\,[P(p,x)\to h(x)\le H],\\
 \mathsf{ListOK}(p,L)&:=
   \bigl(\forall i<|L|\;P(p,L_i)\bigr)\wedge
   \bigl(\forall x\;P(p,x)\to\exists i<|L|\;x=L_i\bigr).
\end{align*}
The two constructors are:
\begin{align*}
 \mathsf{enumerate}(p,H)&:\mathsf{BoundOK}(p,H)	o
     \mathsf{ListOK}(p,\mathsf{enumerate}(p,H)),\\
 \mathsf{maxHeight}(L)&:\mathsf{ListOK}(p,L)	o
     \mathsf{BoundOK}(p,\mathsf{maxHeight}(L)).
\end{align*}
The first depends on a primitive-recursive finite bound on normalized codes of
height at most \(H\), which must be stated for the chosen coding.

\subsection{Uniform packet checklist}

Before a theorem packet is called embedded, record all fields below.

\begin{longtable}{@{}p{0.28\textwidth}p{0.64\textwidth}@{}}
\toprule
Field & Required evidence\\
\midrule
\endhead
Pinned source & ProveIt commit and exact imported module paths.\\
Statement class & Exact, approximate, certificate-relative, relational, or
obstruction theorem.\\
Parameter codes & Natural-number coding and domain guard.\\
Fixed evaluators & Exact \(\PR\) symbols; no unmentioned program input.\\
Challenges/answers & Bounds, verifier, and observational meaning.\\
Compatibility & Adjacent or finite-diagram relation, if outputs are relational.\\
Precision ledger & All input queries, shifts, truncations, margins, and errors.\\
Trace & Initial state, step, bound, invariant, and acceptance predicate.\\
Local proofs & Actual \(\Prf\,\PRAAx\) derivations of base, step, and final
contracts.\\
Search data & Net, height bound, isolating region, cover, or finite candidate
bound where required.\\
Producer status & Proved, conditional, externally certified, or open.\\
Object endpoint & Exact \code{Fm PTm} conclusion and context.\\
Native export & Lean theorem applying \code{pra\_iff\_native} to the completed
endpoint.\\
Audit & Build log, axiom report, placeholder scan, and theorem ledger.\\
Executable checks & Exact examples and negative tests, explicitly labeled as
regression evidence.\\
\bottomrule
\end{longtable}

A packet failing one field may still contain valuable mathematics.  The
checklist identifies the missing layer without conflating it with the layers
already complete.
